\chapter{Home Medicines Review}

\section*{Definition and Overview}
A Home Medicines Review (HMR) is a free, comprehensive review of all the medicines a person is taking, conducted in their own home by an accredited clinical pharmacist in partnership with their GP. The review looks at prescription medicines, over-the-counter products, vitamins, and complementary medicines to identify problems such as interactions, duplication, incorrect dosing, and difficulty taking medicines. The pharmacist provides the GP with recommendations, and the GP and patient agree on any changes. HMRs are available to people in the community who are at risk of medicine-related problems, particularly older people taking multiple medicines.

\section*{Epidemiology}
Medicine-related problems are common and costly. Around one in ten hospital admissions in older Australians is medicine-related, and many are preventable. Older people taking five or more medicines (polypharmacy) are at the highest risk of side effects, interactions, and errors. About half of Australians over 70 take five or more medicines regularly. HMRs have been shown to identify significant medicine problems in the majority of reviews and to reduce hospital admissions and adverse events when acted upon.

\section*{Aetiology and Pathophysiology}
\begin{itemize}
    \item \textbf{Polypharmacy:} multiple medicines increase the risk of interactions and side effects
    \item \textbf{Ageing:} changes in the liver, kidneys, and body composition alter how medicines work
    \item \textbf{Interactions:} prescription medicines can interact with each other, with over-the-counter products, and with complementary medicines
    \item \textbf{Adherence problems:} complex regimens, confusion, and side effects lead to missed or incorrect doses
    \item \textbf{Duplication:} the same or similar medicines may be prescribed by different doctors
    \item \textbf{Inappropriate prescribing:} medicines may no longer be needed or may be outdated
    \item \textbf{Pathophysiology of harm:} unrecognised interactions and side effects cause falls, bleeding, kidney injury, and hospitalisation
\end{itemize}

\section*{Clinical Features}
\begin{itemize}
    \item Taking five or more medicines regularly
    \item Recent hospital admission or a fall
    \item Confusion about which medicines to take and when
    \item Difficulty managing medicines: poor sight, arthritis, or cognitive problems
    \item New symptoms that may be side effects: dizziness, drowsiness, nausea, or falls
    \item Seeing multiple prescribers and using multiple pharmacies
    \item Taking over-the-counter and complementary products alongside prescriptions
\end{itemize}

\section*{Differential Diagnosis}
\begin{itemize}
    \item The review distinguishes medicine-related symptoms from symptoms of the underlying illness
    \item Interactions are identified against known interaction databases
    \item Adherence problems are distinguished from therapeutic failure
    \item Over-the-counter and complementary medicine effects are considered
    \item The GP's clinical knowledge combines with the pharmacist's findings
\end{itemize}

\section*{When to Seek Medical Care}
Anyone at risk of medicine-related problems should ask their GP about an HMR, which is free under Medicare for eligible people. Urgent medical care is needed for serious side effects or suspected overdose.

\textbf{Call 000 (or local emergency services) for:}
\begin{itemize}
    \item Severe drowsiness, confusion, or collapse, which may indicate medicine toxicity
    \item Difficulty breathing or swelling after a new medicine
    \item Bleeding, including nosebleeds, bruising, or blood in stools, on blood thinners
    \item Chest pain or palpitations
    \item Suspected overdose, taking more medicine than prescribed, or another person's medicine
\end{itemize}

\section*{Diagnosis and Investigations}
\begin{itemize}
    \item \textbf{Medicine list:} the patient gathers all medicines, including prescription, over-the-counter, and complementary products
    \item \textbf{The home visit:} the pharmacist reviews storage, dosing, and administration in the home setting
    \item \textbf{Patient interview:} how medicines are actually taken, and any symptoms
    \item \textbf{GP referral:} the GP identifies the patient and provides clinical information
    \item \textbf{The pharmacist's report:} a written assessment with recommendations
    \item \textbf{GP review and action plan:} agreed changes are documented and followed up
    \item \textbf{Medicines list update:} a current list is provided to the patient
\end{itemize}

\section*{Management}
\begin{itemize}
    \item \textbf{Identify and fix problems:} stopping unnecessary medicines, adjusting doses, and simplifying regimens
    \item \textbf{Resolve interactions:} separating doses, changing medicines, or stopping interacting products
    \item \textbf{Improve adherence:} blister packs, dosette boxes, pill organisers, and reminders
    \item \textbf{De-prescribing:} safely reducing medicines that are no longer needed
    \item \textbf{Education:} the patient learns what each medicine is for and how to take it safely
    \item \textbf{Community support:} liaison with community nurses and pharmacists
    \item \textbf{Follow-up:} regular re-reviews when medicines change
\end{itemize}

\section*{Complications}
\begin{itemize}
    \item Falls, fractures, and hospitalisation from sedative and blood pressure medicines
    \item Kidney injury and bleeding from unrecognised interactions
    \item Overdose and poisoning from duplication or confusion
    \item Worsening of the underlying condition from non-adherence
    \item Cost and waste from unused medicines
    \item Delayed recognition of side effects
\end{itemize}

\section*{Prognosis and Outcomes}
HMRs improve safety and quality of life. Studies show they identify clinically significant medicine problems in the majority of patients, reduce hospital admissions, and improve adherence and satisfaction. Most patients have their medicine regimens simplified, with fewer tablets and fewer side effects. The benefits are greatest in older people with multiple medicines, and regular reviews keep the benefits as medicines change over time.

\section*{Prevention and Self-Care}
\begin{itemize}
    \item Ask your GP about a Home Medicines Review if you take several medicines
    \item Keep an up-to-date list of all medicines, including over-the-counter and herbal products
    \item Use one pharmacy so medicines can be checked together
    \item Bring all medicines to every GP appointment
    \item Report new symptoms that may be side effects
    \item Use pill organisers or blister packs if regimens are complex
    \item Ask about stopping medicines that may no longer be needed
    \item Dispose of old medicines safely at a pharmacy
\end{itemize}

\section*{Sources and Further Reading}
The following reputable sources provide further information for healthcare students and patients seeking reliable, current advice about home medicines reviews.

\begin{itemize}
    \item NPS MedicineWise. \textit{Home Medicines Review explained.} https://www.nps.org.au/
    \item Pharmaceutical Society of Australia. \textit{Accredited pharmacists and HMRs.} https://www.psa.org.au/
    \item Healthdirect. \textit{Medicines safety and reviews.} https://www.healthdirect.gov.au/medicines
    \item Australian Government Department of Health. \textit{Medication management programs.} https://www.health.gov.au/
    \item NHS. \textit{Reviews of your medicines.} https://www.nhs.uk/
    \item National Prescribing Service. \textit{Medicine safety information for consumers.} https://www.nps.org.au/
\end{itemize}
